\section{Burnside 与 Pólya：从轨道数到着色计数}
\label{sec:04-Discrete-Mathematics/06-Groups-and-Symmetry/03}

一个圆环上串四颗珠子，每颗黑或白。作坊要开模具，问题是：需要开多少套？两种排法只要转一下能重合，戴在手上就看不出差别，同一套模具就够了。

一共 \(2^4=16\) 种排法。挨个画出来，再拿眼睛两两比对去重——四颗尚可，八颗就是 256 种，比对次数上千，人做不动，程序也未必好写。

上一节已经把``算同一种''这件事写清楚了：全部着色构成集合 \(X\)，允许的动作构成群 \(G\)，答案是轨道的条数 \(\abs{X/G}\)。可轨道并不好数。那张三行的图刚刚展示过，轨道的大小是参差不齐的，1、2、4 都出现了，不能拿 \(\abs{X}\) 除以 \(\abs{G}\) 蒙混过去。

\subsection{换一笔账：数被固定的对象}

难数的是轨道，好数的是另一样东西。固定一个操作 \(g\)，问它把哪些着色原封不动地留下来：

\[
\operatorname{Fix}(g)=\set{x\in X:g\cdot x=x}.
\]

这个集合几乎不用思考就能数。转 \(90^\circ\) 要求四个位置串成一个循环，因而必须同色，只有全黑和全白两种；转 \(180^\circ\) 把位置两两配对，两对各自同色即可，\(2^2=4\) 种；什么都不做则全部 16 种都被固定。

Burnside 引理说，这些好数的量取一次平均，就是要的轨道数。

\begin{theorem}[Burnside]
有限群 \(G\) 作用在有限集 \(X\) 上，则
\[
\abs{X/G}=\frac1{\abs{G}}\sum_{g\in G}\abs{\operatorname{Fix}(g)}.
\]
\end{theorem}

证明是一次数对子。把满足 \(g\cdot x=x\) 的全部数对 \((g,x)\) 拢在一起，按两个方向清点。按 \(g\) 分组，每组的大小是 \(\abs{\operatorname{Fix}(g)}\)，加起来就是等式右边的分子。按 \(x\) 分组，每组的大小是稳定子 \(\abs{G_x}\)；同一条轨道上的点，由轨道--稳定子定理知道稳定子一样大，于是一条轨道整体贡献 \(\abs{G\cdot x}\cdot\abs{G_x}=\abs{G}\)。两种数法结果相同，故对子总数既是 \(\sum_g\abs{\operatorname{Fix}(g)}\)，又是 \(\abs{G}\) 乘以轨道条数。除一下即得。

\begin{center}
\begin{tikzpicture}[x=1cm,y=1cm,font=\footnotesize]
  \foreach \cx in {0,0.67,1.09,1.51,1.93,2.6,3.02,3.44,3.86,4.53,4.95,5.62,6.04,6.46,6.88,7.55}
    \foreach \cy in {0,0.44,0.88,1.32}
      \draw[black!35] (\cx,\cy) rectangle +(0.34,0.34);
  \foreach \cx in {0,0.67,1.09,1.51,1.93,2.6,3.02,3.44,3.86,4.53,4.95,5.62,6.04,6.46,6.88,7.55}
    \fill[blue!45] (\cx,1.32) rectangle +(0.34,0.34);
  \foreach \cx in {0,7.55} {
    \fill[blue!45] (\cx,0.88) rectangle +(0.34,0.34);
    \fill[blue!45] (\cx,0) rectangle +(0.34,0.34);
  }
  \foreach \cx in {0,4.53,4.95,7.55}
    \fill[blue!45] (\cx,0.44) rectangle +(0.34,0.34);
  \foreach \cy/\lab/\num in {1.32/{e}/{16},0.88/{r}/{2},0.44/{r^2}/{4},0/{r^3}/{2}} {
    \node[left] at (-0.14,\cy+0.17) {\(\lab\)};
    \node[right,black!75] at (8.05,\cy+0.17) {\(\num\)};
  }
  \foreach \xa/\xb in {0/0.34, 0.67/2.27, 2.6/4.2, 4.53/5.29, 5.62/7.22, 7.55/7.89} {
    \draw[black!60] (\xa,-0.18) -- (\xb,-0.18);
    \draw[black!60] (\xa,-0.18) -- (\xa,-0.08);
    \draw[black!60] (\xb,-0.18) -- (\xb,-0.08);
    \node[below,black!75] at (0.5*\xa+0.5*\xb,-0.2) {\(4\)};
  }
  \node[right,black!70] at (8.05,-0.05) {\(\;=24\)};
  \node[below,black!70] at (3.9,-0.72) {六条轨道，每条的列和都是 \(4\)，合计 \(6\times4=24\)};
\end{tikzpicture}

\emph{四行是 \(C_4\) 的四个旋转，十六列是全部着色，按轨道分成六组。格子涂色表示这个旋转固定了这个着色。按行加得到 \(16+2+4+2=24\)；按列加，每条轨道无论长短都贡献 \(4=\abs{G}\)，六条轨道同样得到 \(24\)。同一批格子数了两遍，Burnside 引理就是这两遍的对账。}
\end{center}

对四颗珠子代进去，\((16+2+4+2)/4=6\)。作坊开六套模具。

\subsection{循环结构让固定点数可以机械读出}

上面数 \(\operatorname{Fix}\) 时靠的是同一句话：一个位置置换要固定某个着色，被它串在同一个循环上的位置必须同色。于是循环个数一确定，固定点数就出来了——\(q\) 种颜色时

\[
\abs{\operatorname{Fix}(g)}=q^{c(g)},
\]

\(c(g)\) 是 \(g\) 作为位置置换的循环个数。转 \(90^\circ\) 是一个长为 4 的循环，\(c=1\)；转 \(180^\circ\) 是两个长为 2 的循环，\(c=2\)；恒等是四个长为 1 的循环，\(c=4\)。

Pólya 的做法是把整个群的循环结构一次性汇总成一个多项式，称为循环指标：

\[
Z_G(a_1,a_2,\ldots)=\frac1{\abs{G}}\sum_{g\in G}\prod_{k\ge1}a_k^{c_k(g)},
\]

其中 \(c_k(g)\) 是长度为 \(k\) 的循环个数。四位置旋转群的循环指标是

\[
Z_{C_4}=\tfrac14\left(a_1^4+a_2^2+2a_4\right),
\]

把所有 \(a_k\) 都代成颜色数 \(q\)，得到的就是 Burnside 那个平均；\(q=2\) 时回到 6。循环指标多出来的本事在于，把 \(a_k\) 换成颜色权重的 \(k\) 次幂之和（比如两色时换成 \(b^k+w^k\)），展开后每一项的系数直接告诉你``恰好用 \(i\) 颗黑珠''的方案有几种，不必重算。

\subsection{群选错了，答案会精确地回答另一个问题}

上面全程假定手链只能转，不能翻面。如果允许翻过来戴，作用群要换成二面体群 \(D_4\)。四条对称轴给出的置换分两类：穿过两颗珠子的两条轴各留下两个不动位置加一个对换，\(c=3\)；穿过两条边中点的两条轴给出两个对换，\(c=2\)。于是翻折部分贡献 \(8+8+4+4=24\)，加上旋转的 24，平均值 \(48/8=6\)。

答案还是 6。这是四颗珠子的巧合，别当规律。换成六颗珠子两种颜色，只许转的项链有 14 种，允许翻面的手链只有 13 种——有一对项链互为镜像，翻面之后合并了。

该花心思的地方不在平均那一步，而在更早：把哪些动作放进 \(G\)。镜像算不算同一种？珠子的位置有没有编号（有编号就等于什么都不许动，群退化成 \(\set{e}\)，答案回到 16）？颜色的库存是否有限？群列错了，后面的求和即使一位不差，得到的也是另一个问题的正确答案。

还有一层限制值得先说明白：Burnside 只报数量，不交出代表元。它告诉你有六种不同的手链，却不会告诉你是哪六种。实际去重时通常还要另配一件工具——给每条轨道指定一个规范形式（canonical form），例如取该轨道中字典序最小的那个串，然后按规范形式去重。计数与去重是两件事，前者有闭式公式，后者往往得靠搜索。

\subsection{下一步：从计数转向计算}

这一节把对称性用在了数数上：商掉群作用之后还剩多少种构型。同一套语言还有一个用途更广的方向。若不再问``有多少条轨道''，而是问``什么样的函数在整条轨道上取同一个值''，得到的就是不变量；若进一步问``输入沿着轨道移动时，输出应当怎样跟着动''，得到的是等变映射。

下一节沿这个方向走下去。届时会看到，池化层在做的是把整条轨道压成一个数，卷积层在做的是让输出随输入一起平移——这一章的抽象工具在那里变成了具体的网络结构约束。
